\section{注意力的数学形式化}

\subsection{Dot-Product Attention 的数学表示}

\begin{figure}[H]
\centering
\includegraphics[width=0.95\textwidth]{slides-images/slide-021.jpg}
\caption{注意力机制的数学形式化：四个方程完整描述了注意力计算过程\protect\footnotemark}
\end{figure}
\footnotetext{Slides 第22页。}

设编码器隐藏状态为 $\mathbf{h}_1, \ldots, \mathbf{h}_N \in \mathbb{R}^h$，解码器在时间步 $t$ 的隐藏状态为 $\mathbf{s}_t \in \mathbb{R}^h$。注意力计算的完整公式如下：

\textbf{第一步：计算注意力分数}

$$\mathbf{e}^t = [\mathbf{s}_t^T \mathbf{h}_1, \ldots, \mathbf{s}_t^T \mathbf{h}_N] \in \mathbb{R}^N$$

\begin{itemize}
    \item $\mathbf{e}^t$：时间步 $t$ 的注意力分数向量
    \item $\mathbf{s}_t^T \mathbf{h}_i$：解码器状态与编码器第 $i$ 个位置状态的点积
\end{itemize}

\textbf{第二步：通过 Softmax 得到注意力分布}

$$\boldsymbol{\alpha}^t = \mathrm{softmax}(\mathbf{e}^t) \in \mathbb{R}^N$$

\begin{itemize}
    \item $\boldsymbol{\alpha}^t$ 是一个概率分布，所有元素非负且和为1
\end{itemize}

\textbf{第三步：计算注意力输出（上下文向量）}

$$\mathbf{a}_t = \sum_{i=1}^{N} \alpha_i^t \mathbf{h}_i \in \mathbb{R}^h$$

\begin{itemize}
    \item $\mathbf{a}_t$：编码器隐藏状态的加权平均，权重由注意力分布决定
\end{itemize}

\textbf{第四步：拼接并生成}

$$[\mathbf{a}_t; \mathbf{s}_t] \in \mathbb{R}^{2h}$$

将注意力输出与解码器隐藏状态拼接，得到双倍长度的向量，通过线性变换和 softmax 生成输出词的概率分布。

\begin{importantbox}{注意力机制的本质}
注意力机制本质上是一种\textbf{软性信息检索}：给定一个查询（query，即解码器状态），在一组值（values，即编码器状态）中寻找最相关的信息，并通过加权平均的方式提取出来。这种``查询-键值''的思想后来在 Transformer 中被进一步发展为 Query-Key-Value 框架。
\end{importantbox}

\subsection{注意力变体}

注意力机制的通用框架包含一组值（values）$\mathbf{h}_1, \ldots, \mathbf{h}_N \in \mathbb{R}^{d_1}$ 和一个查询（query）$\mathbf{s} \in \mathbb{R}^{d_2}$。不同的注意力变体在``如何计算注意力分数''这一步有所不同。

\begin{figure}[H]
\centering
\includegraphics[width=0.95\textwidth]{slides-images/slide-023.jpg}
\caption{注意力的通用框架：三个核心步骤不变（计算分数、softmax、加权求和），变化的是如何计算注意力分数\protect\footnotemark}
\end{figure}
\footnotetext{Slides 第24页。}

\subsubsection{基本点积注意力（Basic Dot-Product Attention）}

$$e_i = \mathbf{s}^T \mathbf{h}_i \in \mathbb{R}$$

\begin{itemize}
    \item 最简单的形式，直接计算查询和值的点积
    \item 要求 $d_1 = d_2$（查询和值的维度必须相同）
    \item 优点：计算高效，无需额外参数
    \item 缺点：隐藏状态必须同时存储多种信息（输出当前词、记录语法结构、规划未来），并非所有维度都与``回头查看''相关
\end{itemize}

\subsubsection{乘法注意力 / 双线性注意力（Multiplicative / Bilinear Attention）}

$$e_i = \mathbf{s}^T \mathbf{W} \mathbf{h}_i \in \mathbb{R}$$

其中 $\mathbf{W} \in \mathbb{R}^{d_2 \times d_1}$ 是可学习的权重矩阵。

\begin{figure}[H]
\centering
\includegraphics[width=0.95\textwidth]{slides-images/slide-024.jpg}
\caption{注意力变体：基本点积注意力（上）直接做点积；乘法注意力（下）在中间插入一个可学习的矩阵 $\mathbf{W}$\protect\footnotemark}
\end{figure}
\footnotetext{Slides 第25页。参考 Luong, Pham, and Manning 2015。}

\begin{knowledgebox}{为什么乘法注意力更好}
Manning 教授解释道：LSTM 的隐藏状态是它的``完整记忆''，需要存储多种不同类型的信息。编码器和解码器可能将相同类型的信息存储在不同的维度上。乘法注意力通过可学习的矩阵 $\mathbf{W}$，能够学习到``解码器的哪些维度应该与编码器的哪些维度匹配''，而不要求维度对维度的精确对应。这是 Luong, Pham 和 Manning 在2015年提出的方法。
\end{knowledgebox}

\subsubsection{降秩乘法注意力（Reduced-Rank Multiplicative Attention）}

乘法注意力的一个问题是矩阵 $\mathbf{W}$ 的参数量为 $d_1 \times d_2$（如果隐藏状态维度为1000，则有100万个参数）。解决办法是将 $\mathbf{W}$ 分解为两个低秩矩阵：

$$e_i = \mathbf{s}^T (\mathbf{U}^T \mathbf{V}) \mathbf{h}_i = (\mathbf{U}\mathbf{s})^T (\mathbf{V}\mathbf{h}_i)$$

其中 $\mathbf{U} \in \mathbb{R}^{k \times d_2}$，$\mathbf{V} \in \mathbb{R}^{k \times d_1}$，$k \ll d_1, d_2$。

\begin{figure}[H]
\centering
\includegraphics[width=0.95\textwidth]{slides-images/slide-025.jpg}
\caption{降秩乘法注意力（上）：将大矩阵分解为两个小矩阵的乘积，等价于分别对查询和值做低维投影后取点积。加法注意力（下）：使用小型神经网络计算注意力分数\protect\footnotemark}
\end{figure}
\footnotetext{Slides 第26页。}

\begin{importantbox}{从降秩注意力到 Transformer}
降秩乘法注意力的本质是：先将查询和值\textbf{投影到低维空间}，再在低维空间中做点积。这正是 \textbf{Transformer 中多头注意力的做法}！Transformer 取每个大向量，通过投影矩阵将其映射到低维空间，然后在低维空间中计算点积注意力。理解这一点是理解 Transformer 的关键。
\end{importantbox}

\subsubsection{加法注意力（Additive Attention）}

$$e_i = \mathbf{v}^T \tanh(\mathbf{W}_1 \mathbf{h}_i + \mathbf{W}_2 \mathbf{s}) \in \mathbb{R}$$

其中 $\mathbf{W}_1 \in \mathbb{R}^{d_3 \times d_1}$，$\mathbf{W}_2 \in \mathbb{R}^{d_3 \times d_2}$，$\mathbf{v} \in \mathbb{R}^{d_3}$，$d_3$ 是注意力维度（超参数）。

\begin{knowledgebox}{加法注意力的历史地位}
加法注意力是 Bahdanau, Cho, and Bengio（2014）提出的\textbf{最早的注意力形式}。它本质上是一个小型前馈神经网络，用于计算注意力分数。虽然后续研究表明，经过良好的超参数调优后加法注意力可能优于乘法注意力，但由于乘法注意力\textbf{更简单、更高效}，在实践中（特别是 Transformer 中）乘法注意力几乎完全胜出。
\end{knowledgebox}

\begin{table}[H]
\centering
\begin{tabular}{lccc}
\toprule
\textbf{注意力类型} & \textbf{公式} & \textbf{额外参数} & \textbf{提出者} \\
\midrule
基本点积 & $\mathbf{s}^T \mathbf{h}_i$ & 无 & --- \\
乘法/双线性 & $\mathbf{s}^T \mathbf{W} \mathbf{h}_i$ & $\mathbf{W}$ & Luong et al. 2015 \\
降秩乘法 & $(\mathbf{U}\mathbf{s})^T(\mathbf{V}\mathbf{h}_i)$ & $\mathbf{U}, \mathbf{V}$ & (Transformer 使用) \\
加法 & $\mathbf{v}^T \tanh(\mathbf{W}_1 \mathbf{h}_i + \mathbf{W}_2 \mathbf{s})$ & $\mathbf{W}_1, \mathbf{W}_2, \mathbf{v}$ & Bahdanau et al. 2014 \\
\bottomrule
\end{tabular}
\caption{四种注意力变体对比}
\end{table}

\subsection{本章小结}

\begin{itemize}
    \item 注意力机制可以用四个步骤的数学公式完整描述
    \item 核心差异在于如何计算注意力分数：点积、乘法、降秩乘法、加法
    \item 降秩乘法注意力通过低维投影 + 点积实现，这正是 Transformer 的做法
    \item 实践中，乘法/降秩注意力因高效而成为主流
\end{itemize}

%% ========== Part 4: 注意力的优势与通用性 ==========

